\section{Prenex form and Skolemization}
\label{sec:fol-cardinality-normalforms-03}

\begin{definition}[First-order literal and negation normal form]
A \emph{first-order literal} is an atomic formula or the negation of an
atomic formula.  A first-order formula is in \emph{negation normal form} iff
it is generated from first-order literals, \(\bottom\), and
\(\neg\bottom\) by \(\land,\lor,\forall,\exists\).  Thus every negation in
such a formula occurs only in a first-order literal or immediately before
\(\bottom\).  This extends the earlier propositional definition by adding
the two quantifier constructors and supplies the induction class used in the
prenex theorem.
\end{definition}

\begin{definition}[Prenex formula]
A formula is in \emph{prenex form} iff it has the shape
\[
 Q_1x_1\cdots Q_nx_n\,\theta,
\]
where each \(Q_i\) is \(\forall\) or \(\exists\), and \(\theta\) contains no
quantifier.  The formula is \emph{rectified} iff no variable has both a free
and a bound occurrence and distinct quantifiers bind distinct variables.
\end{definition}

\begin{lemma}[First-order negation normal form]
\label{lem:fol-nnf}
There exist functions
\(\operatorname{nnf},\operatorname{nnf}':\Fm(\Lang)\to\Fm(\Lang)\) such
that, for every formula \(\varphi\):
\begin{enumerate}
  \item \(\operatorname{nnf}(\varphi)\) and \(\operatorname{nnf}'(\varphi)\)
        are in negation normal form;
  \item \(\operatorname{nnf}(\varphi)\equiv\varphi\) and
        \(\operatorname{nnf}'(\varphi)\equiv\neg\varphi\); and
  \item \(\FV\) and \(\operatorname{Var}\) of
        \(\operatorname{nnf}(\varphi)\) and of
        \(\operatorname{nnf}'(\varphi)\) equal \(\FV(\varphi)\) and
        \(\operatorname{Var}(\varphi)\).
\end{enumerate}
\textit{Status: complete proof.}
\end{lemma}
\begin{proof}
Define the two functions by simultaneous recursion.  For \(\alpha\)
atomic and \(Q\in\set{\forall,\exists}\),
\[
\begin{aligned}
 \operatorname{nnf}(\alpha)&=\alpha,
 &\operatorname{nnf}'(\alpha)&=\neg\alpha,\\
 \operatorname{nnf}(\bottom)&=\bottom,
 &\operatorname{nnf}'(\bottom)&=\neg\bottom,\\
 \operatorname{nnf}(\neg\varphi)&=\operatorname{nnf}'(\varphi),
 &\operatorname{nnf}'(\neg\varphi)&=\operatorname{nnf}(\varphi),\\
 \operatorname{nnf}(\varphi\land\psi)&=
   \operatorname{nnf}(\varphi)\land\operatorname{nnf}(\psi),
 &\operatorname{nnf}'(\varphi\land\psi)&=
   \operatorname{nnf}'(\varphi)\lor\operatorname{nnf}'(\psi),\\
 \operatorname{nnf}(\varphi\lor\psi)&=
   \operatorname{nnf}(\varphi)\lor\operatorname{nnf}(\psi),
 &\operatorname{nnf}'(\varphi\lor\psi)&=
   \operatorname{nnf}'(\varphi)\land\operatorname{nnf}'(\psi),\\
 \operatorname{nnf}(\varphi\to\psi)&=
   \operatorname{nnf}'(\varphi)\lor\operatorname{nnf}(\psi),
 &\operatorname{nnf}'(\varphi\to\psi)&=
   \operatorname{nnf}(\varphi)\land\operatorname{nnf}'(\psi),\\
 \operatorname{nnf}(Qx\,\varphi)&=Qx\,\operatorname{nnf}(\varphi),
 &\operatorname{nnf}'(\forall x\,\varphi)&=
   \exists x\,\operatorname{nnf}'(\varphi),\\
 &&\operatorname{nnf}'(\exists x\,\varphi)&=
   \forall x\,\operatorname{nnf}'(\varphi).
\end{aligned}
\]
Every recursive call is on an immediate subformula, so the complexity
\(c\) strictly decreases and the recursion is legitimate.

The three claims follow by simultaneous strong induction on
\(c(\varphi)\).  For claim 1, every clause of \(\operatorname{nnf}\) and
\(\operatorname{nnf}'\) produces a first-order literal, \(\bottom\), or
\(\neg\bottom\)---each admitted by the definition of negation normal
form---or combines values of the two functions by
\(\land,\lor,\forall,\exists\); the induction hypotheses close each case.

For claim 2, check each case against the satisfaction clauses.  The
atomic, \(\bottom\), and negation cases are immediate; the binary cases
for \(\operatorname{nnf}'\) are the recursive truth clauses for
\(\neg(\varphi\land\psi)\), \(\neg(\varphi\lor\psi)\), and
\(\neg(\varphi\to\psi)\), and the \(\to\) case of \(\operatorname{nnf}\)
is the truth clause for \(\to\).  The quantifier cases of
\(\operatorname{nnf}'\) are the pointwise computations
\(\Meta,s\entails\neg\forall x\,\varphi\) iff some \(a\in\dom\Meta\) has
\(\Meta,s[x\mapsto a]\entails\neg\varphi\) iff
\(\Meta,s\entails\exists x\,\neg\varphi\), and dually for
\(\exists\)---computed directly from the satisfaction clauses, with the
induction hypothesis applied under each updated assignment.

For claim 3, every clause preserves the multiset of atomic subformulas
and the quantifier occurrences with their variables, so the recursive
\(\FV\) and \(\operatorname{Var}\) clauses compute the same sets on both
sides; the induction hypotheses close each case.
\end{proof}

\begin{lemma}[Rectified variant]
\label{lem:rectified-variant}
For every formula \(\varphi\) and every finite set
\(V\supseteq\operatorname{Var}(\varphi)\) of variables there is a formula
\(r(\varphi,V)\) such that:
\begin{enumerate}
  \item \(r(\varphi,V)\equiv_\alpha\varphi\), hence
        \(r(\varphi,V)\equiv\varphi\);
  \item \(r(\varphi,V)\) is rectified, its binders are pairwise distinct
        and lie outside \(V\), and
        \(\FV(r(\varphi,V))=\FV(\varphi)\);
  \item alpha-equivalent formulas have equal free-variable sets; and
  \item \(r\) preserves the constructor tree; in particular, if
        \(\varphi\) is in negation normal form, then so is
        \(r(\varphi,V)\).
\end{enumerate}
\textit{Status: complete proof.}
\end{lemma}
\begin{proof}
Define \(r\) by recursion: atomic formulas and \(\bottom\) are unchanged;
\(r(\neg\varphi,V)=\neg r(\varphi,V)\);
\(r(\varphi\circ\psi,V)=r(\varphi,V)\circ
r(\psi,\,V\cup\operatorname{Var}(r(\varphi,V)))\) for
\(\circ\in\set{\land,\lor,\to}\); and
\(r(Qx\,\varphi,V)=Qz\,r(\varphi[z/x],\,V\cup\set{z})\), where \(z\) is
the least variable outside \(V\).  The recursion is well founded by the
pattern of the substitution definition: since
\(\operatorname{Var}(\varphi)\subseteq V\) and \(z\notin V\), we have
\(z\notin\operatorname{Var}(\varphi)\cup\set{x}\), so no binder of
\(\varphi\) equals \(z\), the substitution \(\varphi[z/x]\) runs on the
first two quantifier clauses only and is textual replacement---whence
\(\operatorname{Var}(\varphi[z/x])\subseteq
\operatorname{Var}(\varphi)\cup\set{z}\), so the invariant
\(\operatorname{Var}\subseteq V\) persists into every recursive
call---and \(c(\varphi[z/x])=c(\varphi)<c(Qx\,\varphi)\) by
\cref{lem:substitution-complexity}.

Prove claim 3 first.  The relation ``has the same free-variable set''
is an equivalence relation and a congruence for every formula
constructor, by the recursive \(\FV\) clauses.  It contains each
generating pair \(Qx\,\varphi\equiv_\alpha Qy\,(\varphi[y/x])\) with
\(y\notin\operatorname{Var}(\varphi)\): both cases of
\cref{lem:substitution-fv}(2) give
\(\FV(Qy\,(\varphi[y/x]))=\FV(Qx\,\varphi)\), because when
\(x\in\FV(\varphi)\) the sets
\((\FV(\varphi)\setminus\set{x})\cup\set{y}\) minus \(\set{y}\) and
\(\FV(\varphi)\setminus\set{x}\) coincide (\(y\notin\FV(\varphi)\)), and
when \(x\notin\FV(\varphi)\) both sides equal \(\FV(\varphi)\).  By
leastness of \(\equiv_\alpha\), alpha-equivalence is contained in the
relation.

Claims 1, 2, and 4 follow by strong induction on \(c(\varphi)\) under the
invariant \(\operatorname{Var}(\varphi)\subseteq V\).  For claim 1, the
atomic and \(\bottom\) cases are reflexivity, and the connective cases
are the congruence property of \(\equiv_\alpha\).  In the quantifier
case, the induction hypothesis applies to \(\varphi[z/x]\) at complexity
\(c(\varphi[z/x])=c(\varphi)<c(Qx\,\varphi)\)---licensed by
\cref{lem:substitution-complexity}---and gives
\(r(\varphi[z/x],V\cup\set{z})\equiv_\alpha\varphi[z/x]\); the quantifier
congruence gives
\(Qz\,r(\varphi[z/x],V\cup\set{z})\equiv_\alpha Qz\,(\varphi[z/x])\); the
generating clause \(Qx\,\varphi\equiv_\alpha Qz\,(\varphi[z/x])\) is
licensed by \(z\notin\operatorname{Var}(\varphi)\); transitivity and
symmetry conclude.  Then \(r(\varphi,V)\equiv\varphi\) by
\cref{cor:alpha-invariance}.

For claim 2, the binders of \(r(\varphi,V)\) lie outside \(V\) and are
pairwise distinct: the quantifier clause chooses \(z\notin V\) and
recurses with \(V\cup\set{z}\), so all inner binders avoid \(z\) and
\(V\), and the binary clause forbids every variable of the left output to
the right recursion, so the right binders avoid the left binders.  By
claims 1 and 3, \(\FV(r(\varphi,V))=\FV(\varphi)\subseteq
\operatorname{Var}(\varphi)\subseteq V\), so no binder is also free and
\(r(\varphi,V)\) is rectified.

For claim 4, every clause of \(r\) preserves the outer constructor, and
in the quantifier case the textual replacement \(\varphi[z/x]\) changes
only variables inside atomic formulas; induction on the clauses shows
that the constructor tree, and hence membership in the
negation-normal-form class, is preserved.
\end{proof}

\begin{lemma}[Prenex extraction]
\label{lem:prenex-extraction}
For every rectified formula \(\varphi\) in negation normal form there is
a formula \(p(\varphi)\) such that:
\begin{enumerate}
  \item \(p(\varphi)\equiv\varphi\);
  \item \(p(\varphi)\) is rectified and in prenex form, and its matrix is
        quantifier-free and in negation normal form; and
  \item \emph{(invariant)} \(\FV\) and \(\operatorname{Var}\) are
        unchanged, and the prefix of \(p(\varphi)\) consists of exactly
        the quantifier occurrences of \(\varphi\), with the same
        quantifier and the same variable; no new free or bound variables
        occur.
\end{enumerate}
\textit{Status: complete proof.}
\end{lemma}
\begin{proof}
Record two facts.  (C1) If \(y\notin\FV(\chi)\), then
\((Qy\,\pi)\circ\chi\equiv Qy\,(\pi\circ\chi)\) and
\(\chi\circ(Qy\,\pi)\equiv Qy\,(\chi\circ\pi)\) for
\(Q\in\set{\forall,\exists}\) and \(\circ\in\set{\land,\lor}\).  For the
\(\forall\)/\(\lor\) case: if
\(\Meta,s\entails(\forall y\,\pi)\lor\chi\), then either
\(\Meta,s[y\mapsto a]\entails\pi\) for every \(a\), or
\(\Meta,s\entails\chi\) and, by \cref{lem:fol-coincidence},
\(\Meta,s[y\mapsto a]\entails\chi\) for every \(a\); in both cases
\(\Meta,s\entails\forall y\,(\pi\lor\chi)\).  Conversely, if
\(\Meta,s\entails\forall y\,(\pi\lor\chi)\) and
\(\Meta,s\notentails\chi\), then \cref{lem:fol-coincidence} gives
\(\Meta,s[y\mapsto a]\notentails\chi\) for every \(a\), so every
\(s[y\mapsto a]\) satisfies \(\pi\), and
\(\Meta,s\entails\forall y\,\pi\).  The remaining seven cases follow by
the same computation from the satisfaction clauses and
\cref{lem:fol-coincidence}.  (C2) \(\equiv\) is a congruence for
\(\land,\lor\) and for each quantifier, immediately from the
satisfaction clauses.

Define \(p\) by structural induction on the negation-normal generation
of \(\varphi\); every subformula of a rectified formula is rectified, so
the induction hypotheses apply.  On first-order literals, \(\bottom\),
and \(\neg\bottom\), let \(p\) be the identity.  For a quantifier case
\(Qx\,\alpha\), put \(p(Qx\,\alpha)=Qx\,p(\alpha)\): claim 1 is (C2) with
the induction hypothesis, and claims 2 and 3 hold because the prefix of
\(p(Qx\,\alpha)\) consists of \(Qx\) followed by the prefix of
\(p(\alpha)\), whose entries are exactly the quantifier occurrences of
\(\alpha\), pairwise distinct from \(x\) by rectifiedness.

For a binary case \(\alpha\circ\beta\) with
\(\circ\in\set{\land,\lor}\), let
\(p(\alpha)=Q_1y_1\cdots Q_ky_k\,\alpha_0\) and
\(p(\beta)=R_1z_1\cdots R_mz_m\,\beta_0\) with quantifier-free
\(\alpha_0,\beta_0\), and put
\[
 p(\alpha\circ\beta)=
 Q_1y_1\cdots Q_ky_k\,R_1z_1\cdots R_mz_m\,(\alpha_0\circ\beta_0).
\]
The required side conditions hold by rectifiedness of
\(\alpha\circ\beta\) and the invariant: each \(y_i\) is bound in
\(\alpha\), so \(y_i\notin\FV(\beta)=\FV(p(\beta))\); and each
\(z_j\) is bound in \(\beta\), so
\(z_j\notin\FV(\alpha_0)\subseteq\FV(\alpha)\cup\set{y_1,\ldots,y_k}\),
using binder distinctness for \(z_j\ne y_i\).  An inner induction on
\(k+m\) then proves
\(p(\alpha\circ\beta)\equiv p(\alpha)\circ p(\beta)\): while a
quantifier heads either side of \(\circ\), (C1) moves it outward across
\(\circ\)---its variable is not free in the other side by the side
conditions just checked---and (C2) transports the equivalence under the
already-extracted prefix.  With the induction hypotheses
\(p(\alpha)\equiv\alpha\) and \(p(\beta)\equiv\beta\) and (C2), claim 1
follows.  The output is prenex with quantifier-free negation-normal
matrix \(\alpha_0\circ\beta_0\); it is rectified and satisfies the
invariant because nothing is renamed or introduced: its prefix lists
exactly the quantifier occurrences of \(\alpha\) followed by those of
\(\beta\), which are pairwise distinct and disjoint from the free
variables by rectifiedness of \(\alpha\circ\beta\).  The base,
quantifier, and binary cases exhaust the constructors of a first-order
formula in negation normal form.
\end{proof}

\begin{theorem}[Prenex normal form]
\label{thm:prenex-normal-form}
Every first-order formula is logically equivalent to a rectified formula in
prenex form in the same signature.
\textit{Status: complete proof.}
\end{theorem}
\begin{proof}
Let \(\varphi\in\Fm(\Lang)\).  Put
\(\psi=r\bigl(\operatorname{nnf}(\varphi),
\operatorname{Var}(\operatorname{nnf}(\varphi))\bigr)\) and
\(\chi=p(\psi)\), with the three operations supplied by
\cref{lem:fol-nnf,lem:rectified-variant,lem:prenex-extraction}.  By
\cref{lem:fol-nnf}, \(\operatorname{nnf}(\varphi)\equiv\varphi\) and
\(\operatorname{nnf}(\varphi)\) is in negation normal form.  By
\cref{lem:rectified-variant}, \(\psi\) is rectified, remains in negation
normal form, and satisfies \(\psi\equiv\operatorname{nnf}(\varphi)\).  By
\cref{lem:prenex-extraction}, \(\chi\) is rectified and in prenex form
and satisfies \(\chi\equiv\psi\).  Hence
\(\chi\equiv\varphi\), and all three operations stay inside
\(\Fm(\Lang)\), so \(\chi\) is in the same signature.
\end{proof}

\begin{definition}[Skolemization]
Let
\[
 \sigma=Q_1x_1\cdots Q_nx_n\,\theta
\]
be a rectified prenex sentence.  Traverse its quantifier prefix from left to
right.  At an existential quantifier \(\exists x_i\), let
\(u_1,\ldots,u_m\) be the variables governed by preceding universal
quantifiers, introduce a new \(m\)-ary function symbol \(f_i\), replace the
occurrences bound by \(\exists x_i\) with
\(f_i(u_1,\ldots,u_m)\), and delete that existential quantifier.  Retain every
universal quantifier.  The resulting universal sentence in the expanded
signature is the \emph{Skolemization} \(\operatorname{Sk}(\sigma)\).
\end{definition}

An existential quantifier with no preceding universal quantifier produces a
new constant symbol.  A different new function symbol is used for each
existential-quantifier occurrence.

\begin{example}
Skolemization transforms
\[
 \forall x\exists y\forall z\exists w\,R(x,y,z,w)
\]
into
\[
 \forall x\forall z\,R(x,f(x),z,g(x,z)).
\]
The second witness may depend on \(x\) and \(z\); its indirect dependence on
the first witness is already determined by \(f(x)\).
\end{example}

\begin{theorem}[Equisatisfiability of Skolemization]
\label{thm:skolem-equisatisfiable}
Let \(\sigma\) be a rectified prenex \(\Lang\)-sentence.
\begin{enumerate}
  \item \emph{(Axiom of Choice.)}  Every \(\Lang\)-structure satisfying
        \(\sigma\) has an expansion to the Skolem signature satisfying
        \(\operatorname{Sk}(\sigma)\).
  \item \emph{(No choice.)}  The \(\Lang\)-reduct of every model of
        \(\operatorname{Sk}(\sigma)\) satisfies \(\sigma\).
\end{enumerate}
Hence, assuming the Axiom of Choice in the ambient set theory, \(\sigma\)
has a \(\Lang\)-model iff \(\operatorname{Sk}(\sigma)\) has a model in its
expanded signature; the right-to-left implication holds without choice.
\textit{Status: complete proof.}
\end{theorem}
\begin{proof}
Write \(\sigma=Q_1x_1\cdots Q_nx_n\,\theta\) with \(\theta\)
quantifier-free.  For \(0\le i\le n\), let
\(\sigma_i=Q_{i+1}x_{i+1}\cdots Q_nx_n\,\theta\), so that
\(\sigma_0=\sigma\), \(\sigma_n=\theta\), and
\(\FV(\sigma_i)\subseteq\set{x_1,\ldots,x_i}\).  Let \(q(i)\) be the
number of universal positions among the first \(i\), let \(\bar u(j)\)
list the variables of the universal positions before position \(j\), and,
for each existential position \(j\), let \(t_j=f_j(\bar u(j))\) be its
Skolem term.  Then
\(\operatorname{Sk}(\sigma)=\forall u_1\cdots\forall u_p\,\theta^{*}\),
where \(u_1,\ldots,u_p\) are the universal variables in prefix order,
\(p=q(n)\), and \(\theta^{*}=\theta[t_j/x_j]\) simultaneously for all
existential positions \(j\).  Because \(\theta\) is quantifier-free, each
substitution is plain replacement, and because no \(t_j\) contains an
existential variable, the simultaneous substitution equals the
composition of the single substitutions in any order.

Given a structure and a tuple \(\bar a\) of values for the universal
variables among \(x_1,\ldots,x_i\), the \emph{Skolem-term assignment}
\(s_{\bar a}\) assigns to a universal \(x_j\) its entry of \(\bar a\) and
to an existential \(x_j\) the value \(v_j\) of \(t_j\) computed from the
entries for \(\bar u(j)\); the values of all other variables are
irrelevant by \cref{lem:fol-coincidence}.

\emph{Clause 1.}  Let \(\Meta\entails\sigma\).  By finite recursion on
\(i\), choose functions
\(F_j:(\dom\Meta)^{\,q(j-1)}\to\dom\Meta\) for the existential
positions \(j\le i\) while proving the invariant
\[
 \mathrm{Inv}(i):\quad
 \Meta,s_{\bar a}\entails\sigma_i
 \quad\text{for every }\bar a\in(\dom\Meta)^{\,q(i)},
\]
where \(s_{\bar a}\) computes each \(v_j\) as
\(F_j\) applied to the entries of \(\bar a\) for \(\bar u(j)\).
\(\mathrm{Inv}(0)\) is the hypothesis \(\Meta\entails\sigma\).

Universal step (\(Q_{i+1}=\forall\)): from
\(\Meta,s_{\bar a}\entails\forall x_{i+1}\,\sigma_{i+1}\), the universal
satisfaction clause gives
\(\Meta,s_{\bar a}[x_{i+1}\mapsto b]\entails\sigma_{i+1}\) for every
\(b\in\dom\Meta\), which is \(\mathrm{Inv}(i+1)\) for the tuple
\((\bar a,b)\).

Existential step (\(Q_{i+1}=\exists\)): by \(\mathrm{Inv}(i)\) and the
existential satisfaction clause, each witness set
\[
 W(\bar a)=\set{b\in\dom\Meta:
 \Meta,s_{\bar a}[x_{i+1}\mapsto b]\entails\sigma_{i+1}}
\]
is nonempty.  Apply the Axiom of Choice to the set
\(\set{\set{\bar a}\times W(\bar a):\bar a\in(\dom\Meta)^{\,q(i)}}\) of
nonempty sets to obtain a function \(F_{i+1}\) with
\(F_{i+1}(\bar a)\in W(\bar a)\) for every \(\bar a\).  This is one
application of choice per existential quantifier, and choice is used
nowhere else.  With \(v_{i+1}=F_{i+1}(\bar a)\), the Skolem-term
assignment satisfies \(\mathrm{Inv}(i+1)\).

Define the expansion \(\Meta^{*}\) of \(\Meta\) to the Skolem signature
by \(f_j^{\Meta^{*}}=F_j\) for every existential position \(j\).  Fix
\(\bar a\in(\dom\Meta)^{p}\).  By \(\mathrm{Inv}(n)\),
\(\Meta,s_{\bar a}\entails\theta\), hence
\(\Meta^{*},s_{\bar a}\entails\theta\) by \cref{lem:reduct-invariance}.
The term-value clauses give each \(t_j\) the value \(v_j\) under
\([\bar u\mapsto\bar a]\) in \(\Meta^{*}\), so iterating
\cref{thm:fol-substitution} once per existential position turns
\(\Meta^{*},s_{\bar a}\entails\theta\) into
\(\Meta^{*},[\bar u\mapsto\bar a]\entails\theta^{*}\).  Since \(\bar a\)
was arbitrary, the \(p\)-fold universal satisfaction clause yields
\(\Meta^{*}\entails\operatorname{Sk}(\sigma)\).

\emph{Clause 2.}  Let \(\mathcal N\entails\operatorname{Sk}(\sigma)\) in
the expanded signature; now \(v_j\) is the value of the term \(t_j\) in
\(\mathcal N\) at the entries for \(\bar u(j)\), and no choice is needed
to produce it.  By downward induction from \(i=n\) to \(i=0\), prove
\[
 \mathrm{Inv}'(i):\quad
 \mathcal N,s_{\bar a}\entails\sigma_i
 \quad\text{for every }\bar a\in(\dom\mathcal N)^{\,q(i)}.
\]
Base \(i=n\): instantiating the universal prefix of
\(\operatorname{Sk}(\sigma)\) at \(\bar a\) gives
\(\mathcal N,[\bar u\mapsto\bar a]\entails\theta^{*}\), and reversing the
iterated substitution-lemma computation of clause 1, again by
\cref{thm:fol-substitution}, gives
\(\mathcal N,s_{\bar a}\entails\theta\).  Universal step
(\(Q_{i+1}=\forall\)): for every \(b\), \(\mathrm{Inv}'(i+1)\) at
\((\bar a,b)\) gives
\(\mathcal N,s_{\bar a}[x_{i+1}\mapsto b]\entails\sigma_{i+1}\)---the
values \(v_j\) for the existential \(j\le i\) depend only on the entries
of \(\bar a\)---and the universal satisfaction clause gives
\(\mathrm{Inv}'(i)\).  Existential step (\(Q_{i+1}=\exists\)):
\(q(i+1)=q(i)\), and \(\mathrm{Inv}'(i+1)\) at \(\bar a\) exhibits the
witness \(v_{i+1}\), so the existential satisfaction clause gives
\(\mathrm{Inv}'(i)\).  Then \(\mathrm{Inv}'(0)\) says
\(\mathcal N\entails\sigma\) in the expanded signature, and
\cref{lem:reduct-invariance} yields
\(\mathcal N\restr\Lang\entails\sigma\).  No choice principle is used.
\end{proof}

Skolemization generally does not preserve logical equivalence: its output is
in an expanded signature and fixes witness functions that the original
sentence did not name.  The theorem states equisatisfiability and the precise
expansion--reduct relation, not equivalence of formulas.

